\documentclass[11pt,a4paper]{report}
\usepackage[a4paper, margin=2.3cm]{geometry}
\usepackage{fontspec}
\usepackage{xeCJK}
\usepackage[english,spanish,activeacute,es-noshorthands]{babel}
\usepackage{xcolor}
\usepackage{tcolorbox}
\tcbuselibrary{skins, breakable}
\usepackage{booktabs}
\usepackage{longtable}
\usepackage{enumitem}
\usepackage{fancyhdr}
\usepackage{caption}
\captionsetup{font=small, labelfont=bf}
\usepackage{titlesec}
\usepackage{hyperref}
\hypersetup{
  colorlinks=true, linkcolor=blue!50!black, urlcolor=red!50!black,
  pdftitle={Bibliografia AI Espacial - Trilingue ES/EN/ZH},
  pdfauthor={INTISAT}
}

% Fonts
\setmainfont{Latin Modern Roman}
\setCJKmainfont{FandolSong}[
  BoldFont=FandolHei,
  ItalicFont=FandolKai
]

\definecolor{intisat}{RGB}{26, 54, 93}
\definecolor{intisataccent}{RGB}{197, 143, 0}

\titleformat{\chapter}[hang]
  {\normalfont\huge\bfseries\color{intisat}}
  {\thechapter.}{0.6em}{}
\titlespacing*{\chapter}{0pt}{-15pt}{15pt}

\pagestyle{fancy}
\fancyhf{}
\fancyhead[L]{\nouppercase{\leftmark}}
\fancyhead[R]{\thepage}
\renewcommand{\headrulewidth}{0.4pt}

\newtcolorbox{trilingue}[1][]{
  colback=blue!4, colframe=intisat,
  fonttitle=\bfseries, title=Conexion con CubeSat-EdgeAI / Connection with CubeSat-EdgeAI / 与CubeSat-EdgeAI项目的关联,
  breakable, sharp corners=south, #1
}

\begin{document}

% ─── Portada trilingue ──────────────────────────────────────────────
\begin{titlepage}
\centering
\vspace*{1cm}
{\color{intisat}\rule{\textwidth}{2pt}}
\vspace{0.6cm}

{\Huge\bfseries\color{intisat} Inteligencia Artificial en el espacio}\\[0.3em]
{\LARGE\bfseries\color{intisat} Artificial Intelligence in Space}\\[0.3em]
{\LARGE\bfseries\color{intisat} 空间人工智能}

\vspace{0.6cm}
{\color{intisataccent}\rule{0.6\textwidth}{1pt}}
\vspace{0.6cm}

{\Large\itshape Bibliografia comentada (trilingue ES / EN / ZH)\\
Annotated bibliography (trilingual)\\
注释参考文献}

\vspace{1cm}

\begin{tabular}{rl}
\textbf{Proyecto / Project / 项目:} & CubeSat-EdgeAI-Payload \\
\textbf{Plataforma / Platform / 平台:} & INTISAT \\
\textbf{Fecha / Date / 日期:} & \today \\
\textbf{Idiomas / Languages / 语言:} & Español / English / 中文 \\
\end{tabular}

\vfill

\begin{minipage}{0.85\textwidth}
\textbf{Resumen / Summary / 摘要.}\\[0.3em]
\textbf{[ES]} Compilacion trilingue de misiones, hardware, empresas y
papers relevantes para IA a bordo de satelites. Sirve para validar
conceptos y aprender vocabulario tecnico en los tres idiomas.\\[0.3em]
\textbf{[EN]} Trilingual compilation of missions, hardware, companies and
papers relevant to onboard AI in satellites. For concept validation and
learning technical vocabulary across all three languages.\\[0.3em]
\textbf{[ZH]} 关于卫星机载人工智能的相关任务、硬件、公司和论文的三语汇编。
用于概念验证和跨三种语言学习专业词汇。
\end{minipage}

\vspace{1cm}
{\color{intisat}\rule{\textwidth}{2pt}}
\end{titlepage}

\tableofcontents
\clearpage

% =====================================================================
\chapter{Introduccion / Introduction / 介绍}

\noindent\textbf{[ES]}\quad
La inteligencia artificial a bordo de satelites (\emph{onboard AI}) es
una de las areas de mas rapido crecimiento en el sector espacial. Este
documento compila las misiones, empresas, hardware y recursos academicos
mas relevantes, presentados en tres idiomas para permitir al lector
validar el vocabulario tecnico.

\vspace{0.5em}
\noindent\textbf{[EN]}\quad
Onboard AI in satellites is one of the fastest-growing areas in the
space sector. This document compiles the most relevant missions,
companies, hardware and academic resources, presented in three languages
to allow the reader to validate technical vocabulary.

\vspace{0.5em}
\noindent\textbf{[ZH]}\quad
卫星机载人工智能是航天领域增长最快的方向之一。本文档汇编了最重要的
任务、公司、硬件和学术资源，以三种语言呈现，让读者能够在不同语言间
验证专业术语。

\section{Como usar este documento / How to use this document / 如何使用本文档}

\begin{itemize}[leftmargin=2em]
  \item \textbf{[ES]} Cada concepto aparece en los 3 idiomas. Los nombres
        propios (PhiSat, NASA, MarCO) permanecen en su idioma original.
  \item \textbf{[EN]} Each concept appears in all 3 languages. Proper
        nouns remain in their original language.
  \item \textbf{[ZH]} 每个概念都以三种语言出现。专有名词（如PhiSat、NASA、
        MarCO等）保留原文。
\end{itemize}

% =====================================================================
\chapter{Observacion de la Tierra / Earth Observation / 对地观测}

\noindent\textbf{[ES]} El area mas activa de IA a bordo: filtros de
nubes, deteccion de cambios, clasificacion de cultivos, monitoreo de
incendios.

\noindent\textbf{[EN]} The most active area of onboard AI: cloud filters,
change detection, crop classification, wildfire monitoring.

\noindent\textbf{[ZH]} 机载AI最活跃的领域：云层过滤、变化检测、
作物分类、火灾监测。

\begin{center}
\small
\begin{longtable}{p{2.8cm}p{1.2cm}p{3.5cm}p{3cm}p{3cm}}
\toprule
\textbf{Mision / Mission / 任务} &
\textbf{Año / Year / 年份} &
\textbf{Funcion / Function / 功能} &
\textbf{Hardware} &
\textbf{Origen / Origin / 来源} \\
\midrule
\endhead
PhiSat-1 &
2020 &
Filtra nubes con CNN / Cloud filtering CNN / 用CNN过滤云层 &
Intel Movidius Myriad 2 &
ESA \\
PhiSat-2 &
2024 &
6 apps simultaneas / 6 simultaneous apps / 6个同时运行的应用 &
Intel Myriad X &
ESA \\
HYPSO-1/2 &
2022/24 &
ML hyperspectral / Hyperspectral ML / 高光谱机器学习 &
RPi CM4 &
NTNU (Noruega/Norway/挪威) \\
Intuition-1 &
2023 &
IA hyperspectral on-board / Onboard hyperspectral AI / 机载高光谱AI &
Antelope FPGA &
KP Labs (Polonia/Poland/波兰) \\
CogniSAT-6 &
2024 &
EO comercial con AI / Commercial EO with AI / 商业对地观测+AI &
Movidius &
Aiko + Open Cosmos \\
\bottomrule
\end{longtable}
\end{center}

\begin{trilingue}
\textbf{[ES]} Tu INTISAT pertenece a esta familia pero aplicado a
microscopia biologica. PhiSat-1 es el antecedente metodologico mas
cercano.

\textbf{[EN]} Your INTISAT belongs to this family but applied to
biological microscopy. PhiSat-1 is the closest methodological precedent.

\textbf{[ZH]} 你的INTISAT属于这个家族，但应用于生物显微镜领域。
PhiSat-1是方法上最接近的先例。
\end{trilingue}

% =====================================================================
\chapter{Marte y exploracion planetaria / Mars and Planetary Exploration / 火星与行星探索}

\begin{center}
\small
\begin{longtable}{p{2.8cm}p{2.5cm}p{8.5cm}}
\toprule
\textbf{Mision / Mission / 任务} &
\textbf{IA / AI / 人工智能} &
\textbf{Detalle / Detail / 详情} \\
\midrule
\endhead
Curiosity (2012-) &
AEGIS (desde 2016) &
\textbf{[ES]} El rover elige autonomamente que rocas estudiar.\\
&&\textbf{[EN]} The rover autonomously selects which rocks to study.\\
&&\textbf{[ZH]} 火星车自主选择要研究的岩石。\\
\midrule
Perseverance (2021-) &
AutoNav &
\textbf{[ES]} Conduccion autonoma con vision a 5 m/s. Triplica el ritmo vs teleop.\\
&&\textbf{[EN]} Autonomous driving with vision at 5 m/s. Triples pace vs teleop.\\
&&\textbf{[ZH]} 视觉自主驾驶，速度5 m/s。比远程操控快三倍。\\
\midrule
Ingenuity helicopter (2021-2024) &
Vision-only nav &
\textbf{[ES]} SLAM completo a bordo con IMU + camara.\\
&&\textbf{[EN]} Full onboard SLAM with IMU + camera.\\
&&\textbf{[ZH]} 通过IMU和相机实现完整的机载SLAM。\\
\bottomrule
\end{longtable}
\end{center}

% =====================================================================
\chapter{Hardware AI usado en orbita / AI Hardware in Orbit / 在轨AI硬件}

\begin{center}
\small
\begin{longtable}{p{4cm}p{4.5cm}p{6.5cm}}
\toprule
\textbf{Chip} &
\textbf{Donde / Where / 应用于} &
\textbf{Notas / Notes / 备注} \\
\midrule
\endhead
Intel Movidius Myriad 2/X &
PhiSat, CogniSAT &
\textbf{[ES]} El caballo de batalla. $\sim$1 W, 1 TOPS.\\
&&\textbf{[EN]} The workhorse. $\sim$1 W, 1 TOPS.\\
&&\textbf{[ZH]} 业界主力。约1瓦特，1 TOPS算力。\\
\midrule
Google Edge TPU (Coral) &
Tests LEO &
\textbf{[ES]} Validado contra radiacion por NASA y JAXA.\\
&&\textbf{[EN]} Radiation-validated by NASA and JAXA.\\
&&\textbf{[ZH]} 经NASA和JAXA辐射验证。\\
\midrule
NVIDIA Jetson Nano/Orin &
Comercial sats &
\textbf{[ES]} Mas potente pero 5-10 W de consumo.\\
&&\textbf{[EN]} More powerful but 5-10 W consumption.\\
&&\textbf{[ZH]} 更强大，但功耗5-10瓦特。\\
\midrule
AMD Zynq Ultrascale+ &
Mars rovers, JWST &
\textbf{[ES]} Version radiation-hardened. Usado en mision real.\\
&&\textbf{[EN]} Radiation-hardened version. Used in real missions.\\
&&\textbf{[ZH]} 抗辐射加固版本。用于实际任务。\\
\midrule
RPi 4 / CM4 &
HYPSO-1, Astro Pi &
\textbf{[ES]} COTS, no rad-hard. Vida util 1-3 años.\\
&&\textbf{[EN]} COTS, not rad-hard. Lifetime 1-3 years.\\
&&\textbf{[ZH]} 商用现成品，非抗辐射。寿命1-3年。\\
\bottomrule
\end{longtable}
\end{center}

% =====================================================================
\chapter{Empresas principales / Key Companies / 主要公司}

\begin{center}
\small
\begin{longtable}{p{3.5cm}p{2cm}p{8.5cm}}
\toprule
\textbf{Empresa / Company / 公司} &
\textbf{Pais / Country / 国家} &
\textbf{Que hace / What it does / 业务} \\
\midrule
\endhead
ESA $\Phi$-lab &
EU/欧盟 &
\textbf{[ES]} Division de IA de ESA. / \textbf{[EN]} AI division of ESA. / \textbf{[ZH]} 欧空局人工智能部门。 \\
NASA FDL &
USA &
\textbf{[ES]} Sprints anuales de IA aplicada. / \textbf{[EN]} Annual sprints of applied AI. / \textbf{[ZH]} 应用AI年度冲刺项目。 \\
Ubotica &
Irlanda/Ireland/爱尔兰 &
\textbf{[ES]} Hardware AI para PhiSat. / \textbf{[EN]} AI hardware for PhiSat. / \textbf{[ZH]} 为PhiSat提供AI硬件。 \\
KP Labs &
Polonia/Poland/波兰 &
\textbf{[ES]} Procesador edge AI Antelope. / \textbf{[EN]} Antelope edge AI processor. / \textbf{[ZH]} Antelope边缘AI处理器。 \\
Aiko Space &
Italia/Italy/意大利 &
\textbf{[ES]} Software AI para control de mision. / \textbf{[EN]} AI software for mission control. / \textbf{[ZH]} 任务控制AI软件。 \\
Planet Labs &
USA &
\textbf{[ES]} 200+ Doves para observacion terrestre diaria. / \textbf{[EN]} 200+ Doves for daily Earth observation. / \textbf{[ZH]} 200+颗Dove卫星每日对地观测。 \\
Spire Global &
USA &
\textbf{[ES]} Constelacion para meteorologia y AIS. / \textbf{[EN]} Constellation for weather and AIS. / \textbf{[ZH]} 气象与AIS星座。 \\
Satellogic &
Argentina/阿根廷 &
\textbf{[ES]} EO de alta resolucion. / \textbf{[EN]} High-resolution EO. / \textbf{[ZH]} 高分辨率对地观测。 \\
\bottomrule
\end{longtable}
\end{center}

% =====================================================================
\chapter{Frameworks open-source / 开源框架}

\begin{center}
\small
\begin{longtable}{p{3.5cm}p{10cm}}
\toprule
\textbf{Framework} &
\textbf{Funcion / Function / 功能} \\
\midrule
\endhead
NASA cFS &
\textbf{[ES]} Sistema operativo de vuelo con extensiones AI. \\
&\textbf{[EN]} Flight operating system with AI extensions. \\
&\textbf{[ZH]} 带AI扩展的飞行操作系统。 \\
\midrule
NASA F' (FPrime) &
\textbf{[ES]} Framework usado en MarCO e Ingenuity. \\
&\textbf{[EN]} Framework used in MarCO and Ingenuity. \\
&\textbf{[ZH]} 用于MarCO和Ingenuity的框架。 \\
\midrule
ESA Onboard AI demos &
\textbf{[ES]} Repositorio publico de ESA Phi-Lab. \\
&\textbf{[EN]} ESA Phi-Lab public repository. \\
&\textbf{[ZH]} 欧空局Phi-Lab公开代码库。 \\
\midrule
OpenVINO &
\textbf{[ES]} Wrapper para deployment en Movidius. \\
&\textbf{[EN]} Wrapper for Movidius deployment. \\
&\textbf{[ZH]} Movidius部署封装库。 \\
\midrule
TensorFlow Lite Micro &
\textbf{[ES]} Inferencia en microcontroladores con restricciones. \\
&\textbf{[EN]} Inference on constrained microcontrollers. \\
&\textbf{[ZH]} 在资源受限微控制器上进行推理。 \\
\bottomrule
\end{longtable}
\end{center}

% =====================================================================
\chapter{Lectura priorizada / Priority Reading / 优先阅读}

\noindent\textbf{[ES]} Si tu tiempo es limitado, lee en este orden:

\noindent\textbf{[EN]} If your time is limited, read in this order:

\noindent\textbf{[ZH]} 如果时间有限，按此顺序阅读：

\begin{enumerate}[leftmargin=2em]
  \item \textbf{PhiSat-1 paper} (Marcuccio 2021)\\
        \textbf{[ES]} El caso emblematico mas cercano al proyecto.\\
        \textbf{[EN]} The most relevant emblematic case.\\
        \textbf{[ZH]} 与项目最相关的典范案例。

  \item \textbf{Furano et al. AI in Space} (ESA review, 2020)\\
        \textbf{[ES]} El panorama completo en un PDF.\\
        \textbf{[EN]} The complete landscape in one PDF.\\
        \textbf{[ZH]} 全景式综述一份PDF。

  \item \textbf{OBDP 2024 proceedings}\\
        \textbf{[ES]} Lo mas actual del campo.\\
        \textbf{[EN]} The most current in the field.\\
        \textbf{[ZH]} 该领域最新成果。

  \item \textbf{HYPSO papers} (NTNU)\\
        \textbf{[ES]} Usan RPi CM4 con AI, hardware similar al INTISAT.\\
        \textbf{[EN]} Use RPi CM4 with AI, similar hardware to INTISAT.\\
        \textbf{[ZH]} 使用Raspberry Pi CM4+AI，与INTISAT硬件类似。
\end{enumerate}

% =====================================================================
\chapter{Glosario tecnico trilingue / Trilingual Technical Glossary / 三语技术词汇表}

Lista de terminos clave para validar el vocabulario en los tres idiomas.

\begin{center}
\small
\begin{longtable}{p{4cm}p{4cm}p{4.5cm}}
\toprule
\textbf{Español} & \textbf{English} & \textbf{中文} \\
\midrule
\endhead
satelite & satellite & 卫星 \\
nanosatelite & nanosatellite & 纳米卫星 \\
CubeSat & CubeSat & 立方星 \\
orbita & orbit & 轨道 \\
orbita baja terrestre (LEO) & low Earth orbit (LEO) & 低地球轨道 \\
sincronica solar & sun-synchronous & 太阳同步 \\
eclipse & eclipse & 日食/地影 \\
pase (de comunicacion) & pass & 过境 \\
ground station & ground station & 地面站 \\
estacion terrena & ground station & 地面站 \\
carga util & payload & 有效载荷 \\
sensor de imagen & image sensor & 图像传感器 \\
camara & camera & 相机 \\
lensless & lensless & 无透镜 \\
microscopia & microscopy & 显微镜 \\
inteligencia artificial (IA) & artificial intelligence (AI) & 人工智能 \\
aprendizaje profundo & deep learning & 深度学习 \\
red neuronal & neural network & 神经网络 \\
red convolucional (CNN) & convolutional neural network & 卷积神经网络 \\
segmentacion & segmentation & 分割 \\
clasificacion & classification & 分类 \\
modelo & model & 模型 \\
entrenamiento & training & 训练 \\
inferencia & inference & 推理 \\
procesamiento a bordo & onboard processing & 机载处理 \\
edge computing & edge computing & 边缘计算 \\
acelerador AI & AI accelerator & AI加速器 \\
TOPS (tera operaciones / s) & TOPS (tera operations / s) & TOPS（万亿次运算/秒） \\
sistema de potencia (EPS) & power system (EPS) & 电源系统 \\
panel solar & solar panel & 太阳能电池板 \\
bateria & battery & 电池 \\
celda de litio & lithium cell & 锂电池 \\
regulador & regulator & 稳压器 \\
voltaje & voltage & 电压 \\
corriente & current & 电流 \\
potencia & power & 功率 \\
energia & energy & 能量 \\
control de actitud (ADCS) & attitude control (ADCS) & 姿态控制 \\
rueda de reaccion & reaction wheel & 反作用轮 \\
magnetorquer & magnetorquer & 磁力矩器 \\
gyroscopio & gyroscope & 陀螺仪 \\
star tracker & star tracker & 星敏感器 \\
computadora de a bordo (OBC) & on-board computer (OBC) & 星载计算机 \\
software de vuelo & flight software & 飞行软件 \\
telemetria & telemetry & 遥测 \\
comando & command & 指令 \\
uplink & uplink & 上行链路 \\
downlink & downlink & 下行链路 \\
banda UHF & UHF band & UHF频段 \\
banda S & S-band & S频段 \\
modulacion & modulation & 调制 \\
codificacion (FEC) & coding (FEC) & 编码 \\
ancho de banda & bandwidth & 带宽 \\
test ambiental & environmental test & 环境试验 \\
vibracion & vibration & 振动 \\
vacio & vacuum & 真空 \\
termico & thermal & 热 \\
radiacion & radiation & 辐射 \\
sistema de operaciones & operations system & 运行系统 \\
mision & mission & 任务 \\
lanzamiento & launch & 发射 \\
deployment & deployment & 部署 \\
re-entrada & re-entry & 再入 \\
\bottomrule
\end{longtable}
\end{center}

% =====================================================================
\chapter{Frases utiles trilingues / Useful Trilingual Phrases / 实用三语短语}

\noindent\textbf{[ES]} Frases que vas a usar en reuniones tecnicas.

\noindent\textbf{[EN]} Phrases you will use in technical meetings.

\noindent\textbf{[ZH]} 你将在技术会议中使用的短语。

\vspace{1em}

\noindent\textbf{[ES]} ¿Cual es la apertura numerica del sistema?\\
\textbf{[EN]} What is the numerical aperture of the system?\\
\textbf{[ZH]} 系统的数值孔径是多少？

\vspace{0.5em}
\noindent\textbf{[ES]} El consumo en idle es de 3 watts.\\
\textbf{[EN]} Idle consumption is 3 watts.\\
\textbf{[ZH]} 空闲功耗为3瓦特。

\vspace{0.5em}
\noindent\textbf{[ES]} Necesitamos validar el bus I2C.\\
\textbf{[EN]} We need to validate the I2C bus.\\
\textbf{[ZH]} 我们需要验证I2C总线。

\vspace{0.5em}
\noindent\textbf{[ES]} La carga util pesa 800 gramos.\\
\textbf{[EN]} The payload weighs 800 grams.\\
\textbf{[ZH]} 有效载荷重800克。

\vspace{0.5em}
\noindent\textbf{[ES]} La orbita es heliosincrona a 500 kilometros.\\
\textbf{[EN]} The orbit is sun-synchronous at 500 kilometers.\\
\textbf{[ZH]} 轨道是太阳同步轨道，高度500公里。

\vspace{0.5em}
\noindent\textbf{[ES]} El modelo se entrena con 10 mil imagenes.\\
\textbf{[EN]} The model is trained with 10,000 images.\\
\textbf{[ZH]} 该模型使用10000张图像进行训练。

\vspace{0.5em}
\noindent\textbf{[ES]} Hay que aplicar FDIR ante fallo del sensor.\\
\textbf{[EN]} FDIR must be applied if the sensor fails.\\
\textbf{[ZH]} 如果传感器出现故障，必须应用FDIR。

\vspace{0.5em}
\noindent\textbf{[ES]} El downlink usa banda S a 2.2 GHz.\\
\textbf{[EN]} The downlink uses S-band at 2.2 GHz.\\
\textbf{[ZH]} 下行链路使用2.2 GHz的S频段。

\vspace{0.5em}
\noindent\textbf{[ES]} La inferencia tarda 30 segundos por imagen.\\
\textbf{[EN]} Inference takes 30 seconds per image.\\
\textbf{[ZH]} 每张图像的推理需要30秒。

\vspace{0.5em}
\noindent\textbf{[ES]} El bus PC-104 distribuye 5 voltios al payload.\\
\textbf{[EN]} The PC-104 bus distributes 5 volts to the payload.\\
\textbf{[ZH]} PC-104总线向有效载荷分配5伏电压。

\chapter*{Cierre / Closing / 结语}
\addcontentsline{toc}{chapter}{Cierre / Closing / 结语}

\noindent\textbf{[ES]} Este documento es vivo. Cuando descubras un termino
o referencia nueva, agregarla al glosario en los tres idiomas mantiene la
biblioteca util para futuros miembros del equipo.

\vspace{0.5em}
\noindent\textbf{[EN]} This is a living document. When you discover a new
term or reference, adding it to the glossary in all three languages keeps
the library useful for future team members.

\vspace{0.5em}
\noindent\textbf{[ZH]} 这是一份活动文档。当你发现新术语或参考资料时，
将其以三种语言添加到词汇表中，可以为未来的团队成员保持图书馆的实用性。

\vspace{2em}
\noindent
\textbf{Repositorio / Repository / 仓库}:
\url{https://github.com/JaminYC/CubeSat-EdgeAI-Payload}

\end{document}
